
\documentclass[11pt,a4paper]{article}

%% PACHETE MATEMATICE
\usepackage{amsmath,amssymb,amsfonts}
\usepackage{amsthm}
\usepackage{mathpazo}
\usepackage{eulervm}
\usepackage{enumerate}
\usepackage{color}
\usepackage{graphicx}
\usepackage[all]{xy}
\usepackage{xcolor}
\usepackage{framed}
\usepackage[unicode]{hyperref}
\setcounter{MaxMatrixCols}{20}
\usepackage{multicol}
%\usepackage[romanian]{babel}


%% ASPECT
\linespread{1.05}
\usepackage[T1]{fontenc}
\usepackage[utf8x]{inputenc}
\usepackage[margin=2cm]{geometry}
% \usepackage[color]{showkeys}
\usepackage{anyfontsize}
\usepackage{titlesec}
\titleformat*{\section}{\large\bfseries}

\usepackage{fancyhdr}
\pagestyle{fancy}
\rhead{\small \color{gray} Notițe de Adrian Manea}
\lhead{\small \color{gray} BCD-Aero, 2017---2018, L. Costache \& M. Olteanu}


\usepackage[autostyle = false, style = german]{csquotes}
\usepackage[romanian]{babel}
\newcommand\qq{\enquote}

%% COMENZILE MELE
\renewcommand*{\proofname}{Dem.:}
\newcommand\bb{\mathbb}
\newcommand\dr{\mathrm}
\newcommand\kal{\mathcal}
\newcommand\fk{\mathfrak}
\newcommand\op{\oplus}
\newcommand\ot{\otimes}
\newcommand\ds{\displaystyle}
\newcommand\id{\indent}
\newcommand\nid{\noindent}
\newcommand\seq{\subseteq}
\newcommand\sneq{\subsetneq}
\newcommand\speq{\supseteq}
\newcommand\spneq{\supsetneq}
\newcommand\rar{\rightarrow}
\newcommand\Rar{\Rightarrow}
\newcommand\xrar{\xrightarrow}
\newcommand\lar{\leftarrow}
\newcommand\Lar{\Leftarrow}
\newcommand\xlar{\xleftarrow}
\newcommand\lrar{\leftrightarrow}
\newcommand\Lrar{\Leftrightarrow}
\newcommand\ideal{\trianglelefteq}
\newcommand\ai{\text{ a.\^{i}. }}
\newcommand\sk{(\textit{Schi\c{t}\u{a}}) }
\newcommand\rhup{\rightharpoonup}
\newcommand\rhdn{\rightharpoondown}
\newcommand\lhup{\leftharpoonup}
\newcommand\lhdn{\leftharpoondown}
\newcommand\vid{\emptyset}
\newcommand\wh{\widehat}
\newcommand\ol{\overline}
\newcommand\NN{\mathbb{N}}
\newcommand\RR{\mathbb{R}}
\newcommand\ZZ{\mathbb{Z}}
\newcommand\QQ{\mathbb{Q}}
\newcommand\CC{\mathbb{C}}

%% STILURILE MELE
\newtheoremstyle{dotless}{}{}{\itshape}{}{\bfseries}{:}{ }{}

\theoremstyle{dotless}
\newtheorem{theorem}{Teorem\u{a}}[section]
\newtheorem{proposition}{Propozi\c{t}ie}[section]
\newtheorem{lemma}{Lem\u{a}}[section]
\newtheorem{corollary}{Corolar}[section]
\newtheorem{exercise}{Exerci\c{t}iu}

\newtheoremstyle{dotlessdr}{}{}{}{}{\bfseries}{:}{ }{}
\theoremstyle{dotlessdr}
\newtheorem{example}{Exemplu}[section]
\newtheorem{definition}{Defini\c{t}ie}[section]
\newtheorem{remark}{Observa\c{t}ie}[section]




\begin{document}

\begin{center}
  {\large\textbf{Seminar 6 \\ Șiruri și serii de funcții}}
\end{center}

1. Să se arate că șirul de funcții $(f_n)$, unde:
\[
  f_n : \RR \to \RR, \quad f_n(x) = \frac{1}{n} \arctan x^n
\]
converge uniform pe $\RR$, dar:
\[
  \big(\lim_{n \to \infty} f_n(x)\big)^\prime_{x = 1} \neq \lim_{n \to \infty} f'_n(1).
\]

Rezultatele diferă deoarece șirul derivatelor nu converge uniform pe $\RR$.

\vspace{1cm}

2. Să se studieze convergența punctuală și uniformă a șirului de funcții $(f_n)$, cu
\[
  f_n : [0, 1] \to \RR, \quad f_n(x) = nxe^{-nx^2}.
\]

Să se arate că:
\[
  \lim_{n \to \infty} \int_0^1 f_n(x) dx \neq \int_0^1 \lim_{n \to \infty} f_n(x) dx.
\]
Rezultatul se explică prin faptul că șirul nu este uniform convergent.

De exemplu, pentru $x_n = \frac{1}{n} \in [0, 1]$, avem $f_n(x_n) \to 1$, dar
în general $f_n(x) \to 0$ (simplu).

%%%%%%%%%%%%%%%%%%%%%%%%%%%%%%%%%%%%%%%%%%%%%%%%%%%%%%%%%%%%%%%%%%%%%%

\section{Convergența seriilor de funcții}
\label{sec:serii-f}

\begin{definition}\label{def:pc-sf}
  Seria $\sum f_n$ este \emph{simplu (punctual) convergentă către funcția $f$} dacă șirul
  sumelor parțiale $(S_n(x))_n$ este simplu (punctual) convergent către $f$.

  Seria este \emph{uniform convergentă către funcția $f$} dacă șirul sumelor parțiale
  $(S_n(x))_n$ este uniform convergent către $f$.

  Seria $\sum f_n$ este \emph{absolut convergentă} dacă seria $\sum |f_n|$ este
  simplu convergentă.
\end{definition}

Avem următoarele teoreme de derivare și integrare termen cu termen:

\begin{theorem}\label{thm:int-t}
  Fie $\sum f_n$ o serie uniform convergentă de funcții continue $f_n : [a, b] \to \RR$
  și fie $s$ suma acestei serii. Atunci $s$ este o funcție continuă pe $[a, b]$.

  În plus, avem:
  \[
    \int_a^b s(x) dx = \sum_{n \geq 1} \int_a^b f_n(x) dx.
  \]
\end{theorem}

\begin{theorem}\label{thm:der-t}
  Fie $\sum f_n$ o serie punctual convergentă de funcții de clasă $C^1([a, b])$, cu suma
  $s$ pe $[a, b]$ și astfel încît seria derivatelor $\sum f'_n$ să fie uniform convergentă.

  Atunci funcția $s$ este derivabilă pe $[a, b]$ și:
  \[
    s'(x) = \sum_{n \geq 1} f'_n(x), \forall x \in [a, b].
  \]
\end{theorem}

Criteriile de convergență pentru serii de funcții:

\begin{theorem}[Weierstrass]\label{thm:wei-s-f}
  Fie $\sum f_n$, cu $f_n : [a, b] \to \RR$ o serie de funcții și fie $\sum a_n$ o serie
  convergentă de numere reale pozitive.

  Dacă $|f_n(x)| \leq a_n, \forall x \in [a, b]$ și pentru orice $n \geq N$, cu $N$ fixat,
  atunci seria de funcții este uniform convergentă pe $[a, b]$.
\end{theorem}

\begin{definition}\label{def:egal-m}
  Funcțiile $f_n$ se numesc \emph{egal mărginite} dacă există $M \in \RR$, astfel
  încît $f_n \leq M, \forall n \in \NN$.
\end{definition}

Evident, dacă fiecare funcție $f_i$ este mărginită de $m_i$, putem lua
$M = \max_i (m_i)$ și avem egal mărginirea.

\begin{theorem}[Abel]\label{thm:abel-s-f}
  Dacă seria de funcții $\sum f_n$ se poate scrie sub forma $\sum a_n v_n$, astfel
  încît seria de funcții $\sum v_n$ este uniform convergentă, iar $(a_n)$ este
  un șir monoton de funcții egal mărginite, atunci ea este uniform convergentă.
\end{theorem}

\begin{theorem}[Dirichlet]\label{thm:diri-s-f}
  Dacă seria de funcții $\sum f_n$ se poate scrie sub forma $\sum a_n v_n$
  astfel încît șirul sumelor parțiale al seriei $\sum v_n$ să fie un șir de funcții
  egal mărginite, iar $(a_n)_n$ să fie un șir monoton ce converge uniform către 0,
  atunci ea este uniform convergentă.
\end{theorem}

%%%%%%%%%%%%%%%%%%%%%%%%%%%%%%%%%%%%%%%%%%%%%%%%%%%%%%%%%%%%%%%%%%%%%%

\subsection{Exerciții}
\label{subsec:exc-s-f}

3. Să se precizeze convergența seriilor de funcții:
\begin{enumerate}[(a)]
\item $\sum \frac{n^2}{\sqrt{n!}} (x^n + x^{-n}), x \in \Big[\frac{1}{2}, 2\Big]$;
\item $\sum \frac{nx}{1 + n^5 x^2}, x \in \RR$;
\item $\sum \arctan \frac{2x}{x^2 + n^4}, x \in \RR$;
\item $\sum \Big[ e - \Big( 1 + \frac{1}{n} \Big)^n \Big] \cdot \frac{\cos nx}{n + 1}$;
\item $\sum \frac{1}{2^n} \cos (3^n x), x \in \RR$;
\item $\sum x^n(1 - x), x \in [0, 1]$;
\item $\sum \frac{(x + n)^2}{n^4}, x \in [a, b], 0 < a < b$;
\item $\sum \frac{\ln(1 + nx)}{nx^n}, x > 0$;
\end{enumerate}

Indicații:
\begin{enumerate}[(a)]
\item $\frac{1}{2^n} \leq x^n \leq 2^n$, deci seria este $\leq \frac{n^2}{\sqrt{n!}} (2^n + 2^n)$.
  Arătăm acum (cu criteriul raportului) că seria numerică rezultată este convergentă;
\item Găsim maximul funcției (care este în $\sqrt{\frac{1}{n^5}}$), deci seria va fi uniform
  și absolut convergentă, fiind mai mică decît seria convergentă $\sum \frac{1}{2n \sqrt{n}}$;
\item Similar cu raționamentul anterior;
\item
  \[
    \Big| \big[ e - \big(1 + \frac{1}{n}\big)^n \big] \frac{\cos nx}{n+1} \Big| \leq
    \Big[ \big( 1 + \frac{1}{n} \big)^{n + 1} - \big(1 + \frac{1}{n}\big)^n \Big] \cdot \frac{1}{n+1} < \frac{3}{n(n+1)}.
  \]
  Seria numerică rezultată este acum convergentă, putînd fi comparată cu o serie armonică.
\end{enumerate}

\vspace{1cm}

4. Să se studieze convergența seriilor de funcții și să se decidă dacă
se pot deriva termen cu termen (indicație: Weierstrass, comparație cu
seria armonică):
\begin{enumerate}[(a)]
\item $\sum n^{-x}, x \in \RR$;
\item $\sum \frac{\sin nx}{2^n}, x \in \RR$;
\item $\sum \frac{\sin nx}{n(n+1)}, x \in \RR$;
\item $\sum (-1)^n \frac{x^2 + n}{n^2}, x \in \RR$;
\end{enumerate}

\vspace{1cm}

5. Fie seria $\sum \frac{nx^n + x}{n^2 + 1}$.
\begin{enumerate}[(a)]
\item Pentru ce valori ale lui $x \in \RR$ seria converge?
\item Să se studieze convergența uniformă.
\item Se poate deriva seria termen cu termen?
\end{enumerate}

6. Să se stabilească natura seriei $\sum (f_n - f_{n-1})$, dacă:
\begin{enumerate}[(a)]
\item $f_n : [0, 1] \to \RR, f_n(x) = nx(1 - x)^n, \forall n \in \NN$;
\item $f_n : (0, 1] \to \RR, f_n(x) = \frac{e^{nx}}{1 + e^{nx}}, \forall n \in \NN$.
\end{enumerate}

\vspace{1cm}

7. Să se determine mulțimea de convergență pentru următoarele serii de funcții:
\begin{enumerate}[(a)]
\item $\sum \Big( 1 + \frac{1}{n} \Big)^n \Big( \frac{1 - x}{1 - 2x} \Big)^n, x \neq \frac{1}{2}$;
  (rădăcină pentru seria valorilor absolute)
\item $\sum (-1)^n \frac{1}{\ln n} \Big( \frac{1 - x^2}{1 + x^2} \Big)^n, x \in \RR$; (raport)
\item $\sum 2^n \sin \frac{x}{3^n}, x \in \RR$; (comparație cu $\Big( \frac{2}{3} \Big)^n$)
\item $\sum \frac{\ln(1 + a^n)}{n^x}, a \geq 0$; ($0 < a < 1$ raport, $a = 1$ C, $a > 1$ descompunem în două serii, $a = 0$ C)
\item $\sum \frac{\sin^n x}{n^a}, a \in \RR$. (radical)
\end{enumerate}


\end{document}


